\documentclass[10pt,a4paper]{amsart}
\usepackage[margin=19mm]{geometry}
\usepackage{amsmath,amssymb,mathtools}
\usepackage[scr=rsfs,cal=euler]{mathalfa}
\usepackage{xfrac}
\usepackage[unicode,hidelinks,bookmarks=false]{hyperref}
\newcommand{\ie}{i.e.\ }
\newcommand{\cf}{cf.\ }
\newcommand{\resp}{respectively\ }
\newcommand{\Subsection}{\subsection}
\providecommand{\nobreakdash}{\nobreak\hbox{-}}
\setlength{\emergencystretch}{2em}
\multlinegap0pt
\usepackage{amssymb}
\usepackage{xcolor}
\newtheorem{prop}{Proposition}
\newtheorem{thm}{Theorem}
\newcommand{\bR}{\mathbb{R}}
\newcommand{\cB}{\mathcal{B}}
\newcommand{\oM}{\operatorname{M}}
\newcommand{\oP}{\operatorname{P}}
\newcommand{\oS}{\operatorname{S}}
\title{The impossibility of extending\\the Naimark complement}
\author{Emily J. King, Dustin G. Mixon}
\date{Source: arXiv:2504.09534v2 (CC BY 4.0)}
\begin{document}
\maketitle

\noindent\emph{Source and licence.} The impossibility of extending\\the Naimark complement. Version arXiv:2504.09534v2 (CC BY 4.0), distributed by arXiv under the Creative Commons Attribution 4.0 International licence, http://creativecommons.org/licenses/by/4.0/, as recorded in the arXiv licence metadata for this exact version and shown on the abstract page https://arxiv.org/abs/2504.09534v2. Full licence notice: \texttt{licenses/arxiv-2504.09534-CC-BY-4.0.txt}.\par\smallskip

\noindent\textit{Editorial scope.} Counted unit: the article's thm thm.main (source0.tex lines 193--203). The article's complete original proof of it is delivered with it (source0.tex lines 205--301, one \texttt{proof} environment of 97 lines). No mathematics is written, repaired or re-derived: the statement and the proof are the article's own lines, in the article's own notation, and no retained line is substituted or re-flowed.  Retained uncounted as the source-local context the proof needs: lines 144--148; lines 181--186.  Nothing was omitted from any retained span, so the package carries no omission record.  The retained lines cite 2 works, whose entries are transcribed verbatim from the article's own bibliography source in the e-print and delivered with short editorial labels recorded in the item's reference mapping.  No retained line contains a figure environment, an image-inclusion command or drawing code, so no image is delivered and the package declares no asset dependency.  Editorial formatting only, in addition: an AMS \texttt{amsart} preamble that restates the article's own declarations and macros used by the retained lines, namely the environments \texttt{prop}, \texttt{thm} on separate counters, together with the standard packages those lines need.  The retained environments print in this document as Proposition 1, Theorem 1 in order of appearance, whereas the article prints the same items with the numbering of its own full document.  The complete native source of the article as distributed in the arXiv e-print for this exact version is bundled unchanged as \texttt{sources/176448/source0.tex}.
\begin{prop}
\label{prop:matroid}
Consider any matroid $M=(X,\cB)$ that is represented by the column vectors of a rank-deficient nonzero matrix $A\in\bR^{n\times n}$.
The Gale dual $M^\ast$ is represented by the column vectors of the $n \times n$ orthogonal projection onto the kernel of $A$.
\end{prop}

\begin{equation}
\label{eq.casazza}
E\colon\oS(n)\to\oS(n),
\qquad
E(A)=\lambda_{\max}(A)I-A.
\end{equation}

\begin{thm}
\label{thm.main}
For $n\geq2$, the Naimark complement over $\oP(n)$ has an extension over $\oS(n)$ that is
\begin{itemize}
\item[(a)] stratum-wise continuous,
\item[(b)] involutive, and
\item[(c)] Gale
\end{itemize}
if and only if $n=2$.
Otherwise, there is an extension that satisfies any two of (a), (b), and (c).
\end{thm}

\begin{proof}
First, we prove ($\Rightarrow$) by contradiction.
Take any $n>2$, and suppose there exists a stratum-wise continuous involutive Gale extension $E\colon\oS(n)\to\oS(n)$ of the Naimark complement.
Notice that for any $A\in\oS(n)$, the size of every basis in the matroid of $A$ equals the rank of $A$, and so
\[
\operatorname{rank}E(A)
=n-\operatorname{rank}(A)
\]
since $E$ is Gale.
Next, since $E$ is involutive, then for each $d\in\{1,\ldots,n-1\}$, $E$ determines a bijection between the rank-$d$ and rank-$(n-d)$ strata of $\oS(n)$.
Since $E$ is stratum-wise continuous, these bijections are homeomorphisms.
Meanwhile, the dimension of the rank-$d$ stratum is $dn-\binom{d}{2}$ (see~\cite{BonnabelS:10}, for example).
As such, the rank-$1$ and rank-$(n-1)$ strata are homeomorphic manifolds of dimensions $n$ and $n(n-1)-\binom{n-1}{2}>n$, which contradicts invariance of domain.

For the remainder of the proof, we construct extensions of the Naimark complement that satisfy various combinations of (a), (b), and (c).

\medskip
\noindent
\textbf{Case I:} (a) and (b) (and also (c) when $n=2$).
Consider the map $E$ defined in Equation~\eqref{eq.casazza}.
First, for every $P\in\oP(n)$, it holds that
\[
\lambda_{\max}(P)I-P
=I-P
=N(P),
\]
and so $E$ is an extension of $N$.
Next, (a) follows from the fact that $A\mapsto\lambda_{\max}(A)$ is continuous (which in turn is a consequence of Weyl's inequality).
For (b), observe that
\[
\lambda_{\max}(E(A))
=\lambda_{\max}(\lambda_{\max}(A)I-A)
=\lambda_{\max}(A)-\lambda_{\min}(A)
=\lambda_{\max}(A),
\]
and so
\[
E(E(A))
=\lambda_{\max}(E(A))I-E(A)
=\lambda_{\max}(A)I-(\lambda_{\max}(A)I-A)
=A.
\]
In the special case where $n=2$, note that $\oS(2)$ consists of all positive scalar multiples of the members of $\oP(2)$ as a consequence of the real spectral theorem:
\[
\oS(2)
=\{\lambda P:\lambda>0,P\in\oP(2)\}.
\]
In this case, $E$ can be rewritten as
\[
E(\lambda P)=\lambda N(P)=\lambda(I-P).
\]
Since $P$ and $I-P$ have Gale-dual matroids, and since multiplying by $\lambda>0$ does not change these matroids, it follows that $E$ is Gale in this case.

\medskip
\noindent
\textbf{Case II:} (a) and (c).
Consider the map $E$ that sends any matrix to the orthogonal projection onto its kernel.
This is clearly an extension of $N$.
For (a), we apply the Davis--Kahan $\sin\Theta$ theorem~\cite{DavisK:70}.
Consider any $A\in\oS(n)$ of rank $d\in\{1,\ldots,n-1\}$, and let $\lambda>0$ denote its $d$th largest eigenvalue.
Given any $\epsilon>0$, put $\delta:=\lambda\epsilon/\sqrt{2}$.
Then for every $B\in\oS(n)$ of rank $d\in\{1,\ldots,n-1\}$ satisfying $\|A-B\|_F<\delta$, it holds that
\[
\|E(A)-E(B)\|_F
=\sqrt{2}\cdot\|\sin\Theta\|_F
\leq\sqrt{2}\cdot\frac{\|A-B\|_F}{\lambda}
<\epsilon,
\]
where $\sin\Theta$ denotes the diagonal matrix whose diagonal entries are the sines of the principal angles between the kernels of $A$ and of $B$.
Thus, the restriction of $E$ to the rank-$d$ stratum of $\oS(n)$ is continuous at $A$, and (a) follows since $A$ and $d$ were arbitrary.
Meanwhile, (c) follows immediately from Proposition~\ref{prop:matroid}.

\medskip
\noindent
\textbf{Case III:} (b) and (c).
Let $\oM(n)$ denote the (finite) set of matroids that can be represented by members of $\oS(n)$.
Proposition~\ref{prop:matroid} establishes that for each $M\in\oM(n)$, the Gale dual $M^*$ is also a member of $\oM(n)$. 
Let $S_M$ denote the set of matrices in $\oS(n)$ with matroid $M\in\oM(n)$.
Considering $S_M$ is closed under positive scalar multiplication, it holds that
\[
|\mathbb{R}|
=|\mathbb{R}_+\setminus\{1\}|
\leq|S_M\setminus\oP(n)|
\leq|\mathbb{R}^{n\times n}|
=|\mathbb{R}|.
\]
The Schr\"{o}der--Bernstein theorem then delivers a bijection $B_M\colon S_M\setminus{P}(n)\to S_{M^*}\setminus{P}(n)$ for each $M\in\oM(n)$.
Note that we can iteratively redefine $B_{M^*}$ to be the inverse of $B_M$ for each $M\in\oM(n)$.
The resulting bijections together determine a map $E\colon\oS(n)\to\oS(n)$ by
\[
E(A)
=\left\{\begin{array}{rl} N(A)&\text{if }A\in\oP(n)\\
B_M(A)&\text{if }A\in S_M\setminus\oP(n),
\end{array}\right.
\]
which extends $N$ in a way that satisfies (b) and (c).
\end{proof}

\begin{thebibliography}{99}
\bibitem[Bonnab]{BonnabelS:10} S.\ Bonnabel, R.\ Sepulchre, Riemannian metric and geometric mean for positive semidefinite matrices of fixed rank, SIAM J.\ Matrix Anal.\ Appl.\ 31 (2010) 1055--1070.
\bibitem[DavisK]{DavisK:70} C.\ Davis, W.\ M.\ Kahan, The rotation of eigenvectors by a perturbation III, SIAM J.\ Numerical Anal.\ 7 (1970) 1--46.
\end{thebibliography}
\end{document}
